% Relatório da Parte 3 de COS242. Os números vêm de tabelas.tex, gerado por tabelas.py a partir
% de data.json (bench/summarize_p3.py). Compilar: ver README.md.
\documentclass[a4paper,11pt]{article}

\usepackage[brazil]{babel}
\usepackage[T1]{fontenc}
\usepackage{libertinus}
\setmonofont{Inconsolatazi4}[Extension=.otf, UprightFont=*-Regular, BoldFont=*-Bold, Scale=0.92]
\usepackage[a4paper,hmargin=2.2cm,top=1.8cm,bottom=1.9cm]{geometry}
\usepackage{microtype}
\usepackage{amsmath}
\usepackage{graphicx}
\usepackage{booktabs}
\usepackage{array}
\usepackage[table]{xcolor}
\usepackage{listings}
\usepackage{caption}
\usepackage{enumitem}
\usepackage[hidelinks]{hyperref}

% ── tipografia e espaçamento (como na Parte 2, um pouco mais compacto) ──
\linespread{1.12}
\setlength{\parindent}{1.2em}
\setlength{\parskip}{1.5pt plus 1pt}
\setlist{nosep,leftmargin=1.5em,topsep=2pt}
\captionsetup{font=small,labelfont=bf,labelsep=period,skip=4pt}
\captionsetup[table]{position=top}
\setlength{\textfloatsep}{9pt plus 2pt minus 3pt}
\setlength{\floatsep}{8pt plus 2pt minus 2pt}
\setlength{\intextsep}{8pt plus 2pt minus 2pt}
\renewcommand{\arraystretch}{1.08}

\definecolor{heapc}{HTML}{1B7F79}
\definecolor{ruim}{HTML}{C0392B}
\definecolor{codkw}{HTML}{1F4E9C}
\definecolor{codtipo}{HTML}{1B7F79}
\definecolor{codcom}{HTML}{7A7A7A}
\lstset{language=C++, basicstyle=\ttfamily\scriptsize, columns=fullflexible, keepspaces=true,
  keywordstyle=\color{codkw}\bfseries, commentstyle=\color{codcom},
  numbers=none, xleftmargin=1.2em, frame=none, aboveskip=3pt, belowskip=3pt}
\renewcommand{\lstlistingname}{Listagem}
\newcommand{\cod}[1]{\texttt{#1}}
\newcommand{\W}[1]{W\textsubscript{#1}}

\makeatletter
\newcommand{\dado}[1]{\@ifundefined{dado@#1}{\textbf{??\detokenize{#1}??}}{\csname dado@#1\endcsname}}
\renewcommand\section{\@startsection{section}{1}{\z@}{-1.9ex plus -.5ex minus -.2ex}%
  {0.8ex plus .2ex}{\normalfont\large\bfseries}}
\renewcommand\subsection{\@startsection{subsection}{2}{\z@}{-1.5ex plus -.4ex minus -.2ex}%
  {0.4ex plus .2ex}{\normalfont\normalsize\bfseries}}
\makeatother
% Gerado por report/parte3/tabelas.py a partir de data.json; não edite à mão.
\newcommand{\tabelaGrafos}{\W{1} & \dado{W_1.V} & \dado{W_1.E} & \dado{W_1.deg} & \dado{W_1.neg} & \dado{W_1.fwmem}\\
\W{2} & \dado{W_2.V} & \dado{W_2.E} & \dado{W_2.deg} & \dado{W_2.neg} & \dado{W_2.fwmem}\\
\W{3} & \dado{W_3.V} & \dado{W_3.E} & \dado{W_3.deg} & \dado{W_3.neg} & \dado{W_3.fwmem}\\
\W{4} & \dado{W_4.V} & \dado{W_4.E} & \dado{W_4.deg} & \dado{W_4.neg} & \dado{W_4.fwmem}\\
\W{5} & \dado{W_5.V} & \dado{W_5.E} & \dado{W_5.deg} & \dado{W_5.neg} & \dado{W_5.fwmem}\\}
\newcommand{\tabelaDistancias}{10 & $-\infty$ & 10,47 & 4,63 & 9,04 & 37,85\\
20 & $-\infty$ & 8,48 & 5,42 & 8,83 & 42,63\\
30 & $-\infty$ & 8,71 & 4,67 & 9,23 & 37,86\\}
\expandafter\def\csname matriz@W_1\endcsname{\begin{tabular}{@{}r|rrrrr@{}}
 & \textbf{10} & \textbf{20} & \textbf{30} & \textbf{40} & \textbf{50}\\ \midrule
\textbf{10} & --- & --- & --- & --- & ---\\
\textbf{20} & --- & --- & --- & --- & ---\\
\textbf{30} & --- & --- & --- & --- & ---\\
\textbf{40} & --- & --- & --- & --- & ---\\
\textbf{50} & --- & --- & --- & --- & ---\\
\end{tabular}}
\expandafter\def\csname matriz@W_2\endcsname{\begin{tabular}{@{}r|rrrrr@{}}
 & \textbf{10} & \textbf{20} & \textbf{30} & \textbf{40} & \textbf{50}\\ \midrule
\textbf{10} & 0 & 11,73 & 8,14 & 9,49 & 7,08\\
\textbf{20} & 3,84 & 0 & 7,03 & 8,30 & 4,70\\
\textbf{30} & 5,79 & 12,55 & 0 & 8,30 & 6,33\\
\textbf{40} & 5,05 & 12,32 & 5,36 & 0 & 5,71\\
\textbf{50} & 4,48 & 12,23 & 7,79 & 8,26 & 0\\
\end{tabular}}
\expandafter\def\csname matriz@W_3\endcsname{\begin{tabular}{@{}r|rrrrr@{}}
 & \textbf{10} & \textbf{20} & \textbf{30} & \textbf{40} & \textbf{50}\\ \midrule
\textbf{10} & 0 & 4,99 & 5,71 & 7,49 & 4,31\\
\textbf{20} & 3,48 & 0 & 6,19 & 8,27 & 4,32\\
\textbf{30} & 3,96 & 5,80 & 0 & 8,30 & 3,17\\
\textbf{40} & 3,58 & 5,59 & 6,10 & 0 & 3,84\\
\textbf{50} & 3,25 & 5,41 & 6,12 & 8,10 & 0\\
\end{tabular}}
\expandafter\def\csname matriz@W_4\endcsname{\begin{tabular}{@{}r|rrrrr@{}}
 & \textbf{10} & \textbf{20} & \textbf{30} & \textbf{40} & \textbf{50}\\ \midrule
\textbf{10} & 0 & 10,52 & 10,15 & 10,69 & 8,24\\
\textbf{20} & 7,79 & 0 & 9,19 & 9,67 & 7,58\\
\textbf{30} & 7,61 & 10,25 & 0 & 8,77 & 6,55\\
\textbf{40} & 7,94 & 8,65 & 9,22 & 0 & 7,75\\
\textbf{50} & 7,99 & 11,41 & 8,80 & 8,55 & 0\\
\end{tabular}}
\expandafter\def\csname matriz@W_5\endcsname{\begin{tabular}{@{}r|rrrrr@{}}
 & \textbf{10} & \textbf{20} & \textbf{30} & \textbf{40} & \textbf{50}\\ \midrule
\textbf{10} & 0 & 40,20 & 39,56 & 41,55 & 44,91\\
\textbf{20} & 39,17 & 0 & 39,98 & 49,82 & 45,16\\
\textbf{30} & 34,46 & 42,64 & 0 & 47,04 & 43,85\\
\textbf{40} & 38,73 & 43,91 & 47,30 & 0 & 45,24\\
\textbf{50} & 44,06 & 40,62 & 45,25 & 49,11 & 0\\
\end{tabular}}
\newcommand{\tabelaTempos}{\W{1} & \dado{W_1.bf.rounds} & 102 $\pm$ 51\,\textmu s\textsuperscript{$\dagger$} & 602,9 $\pm$ 79,4\,ms\textsuperscript{$\dagger$} & 267 $\pm$ 173\,\textmu s\textsuperscript{$\dagger$}\\
\W{2} & \dado{W_2.bf.rounds} & 12,3 $\pm$ 1,1\,ms & 8,3 $\pm$ 0,7\,min & 38,14 $\pm$ 1,79\,s\\
\W{3} & \dado{W_3.bf.rounds} & 102,7 $\pm$ 4,1\,ms & \textit{memória} & \textit{2,9 $\pm$ 0,1\,h}\\
\W{4} & \dado{W_4.bf.rounds} & 1,47 $\pm$ 0,05\,s & \textit{memória} & \textit{14,4 $\pm$ 0,6\,dias}\\
\W{5} & \dado{W_5.bf.rounds} & 6,11 $\pm$ 0,06\,s & \textit{memória} & \textit{2,0 $\pm$ 0,1\,anos}\\}
\newcommand{\tabelaQuatro}{\W{1} & \dado{W_1.ops.fw} & \dado{W_1.ops.bf} & 602,9\,ms\textsuperscript{$\dagger$} & 267\,\textmu s\textsuperscript{$\dagger$} & ---\\
\W{2} & \dado{W_2.ops.fw} & \dado{W_2.ops.bf} & 8,3\,min & 38,14\,s & 13,1$\times$\\
\W{3} & \dado{W_3.ops.fw} & \dado{W_3.ops.bf} & \textit{8,9\,h} & \textit{2,9\,h} & 33,7$\times$\\
\W{4} & \dado{W_4.ops.fw} & \dado{W_4.ops.bf} & \textit{369\,dias} & \textit{14,4\,dias} & 282$\times$\\
\W{5} & \dado{W_5.ops.fw} & \dado{W_5.ops.bf} & \textit{1.012\,anos} & \textit{2,0\,anos} & 5.443$\times$\\}
\expandafter\def\csname dado@cpu\endcsname{Apple M3 Pro}
\expandafter\def\csname dado@ram\endcsname{36\,GiB}
\expandafter\def\csname dado@threads\endcsname{11}
\expandafter\def\csname dado@compiler\endcsname{Apple clang version 21.0.0}
\expandafter\def\csname dado@limit\endcsname{16\,GiB}
\expandafter\def\csname dado@tests\endcsname{55.823.939}
\expandafter\def\csname dado@W_1.V\endcsname{25.000}
\expandafter\def\csname dado@W_1.E\endcsname{549.953}
\expandafter\def\csname dado@W_1.deg\endcsname{44,0}
\expandafter\def\csname dado@W_1.neg\endcsname{sim}
\expandafter\def\csname dado@W_1.fwmem\endcsname{7,5\,GB}
\expandafter\def\csname dado@W_1.fwmem.gib\endcsname{7,0\,GiB}
\expandafter\def\csname dado@W_1.bf\endcsname{102\,\textmu s}
\expandafter\def\csname dado@W_1.bf.rounds\endcsname{4}
\expandafter\def\csname dado@W_1.bf.runs\endcsname{10}
\expandafter\def\csname dado@W_1.fw\endcsname{602,9\,ms}
\expandafter\def\csname dado@W_1.ap\endcsname{267\,\textmu s}
\expandafter\def\csname dado@W_1.ap.sample\endcsname{25.000}
\expandafter\def\csname dado@W_1.ops.fw\endcsname{$1{,}56\cdot10^{13}$}
\expandafter\def\csname dado@W_1.ops.bf\endcsname{$5{,}50\cdot10^{10}$}
\expandafter\def\csname dado@W_1.ops.ratio\endcsname{284$\times$}
\expandafter\def\csname dado@W_1.ap.threads\endcsname{11}
\expandafter\def\csname dado@W_1.fwest\endcsname{---}
\expandafter\def\csname dado@W_2.V\endcsname{25.000}
\expandafter\def\csname dado@W_2.E\endcsname{824.964}
\expandafter\def\csname dado@W_2.deg\endcsname{66,0}
\expandafter\def\csname dado@W_2.neg\endcsname{não}
\expandafter\def\csname dado@W_2.fwmem\endcsname{7,5\,GB}
\expandafter\def\csname dado@W_2.fwmem.gib\endcsname{7,0\,GiB}
\expandafter\def\csname dado@W_2.bf\endcsname{12,3\,ms}
\expandafter\def\csname dado@W_2.bf.rounds\endcsname{24}
\expandafter\def\csname dado@W_2.bf.runs\endcsname{10}
\expandafter\def\csname dado@W_2.fw\endcsname{8,3\,min}
\expandafter\def\csname dado@W_2.ap\endcsname{38,14\,s}
\expandafter\def\csname dado@W_2.ap.sample\endcsname{25.000}
\expandafter\def\csname dado@W_2.ops.fw\endcsname{$1{,}56\cdot10^{13}$}
\expandafter\def\csname dado@W_2.ops.bf\endcsname{$4{,}95\cdot10^{11}$}
\expandafter\def\csname dado@W_2.ops.ratio\endcsname{31,6$\times$}
\expandafter\def\csname dado@W_2.bf.nsedge\endcsname{0,62}
\expandafter\def\csname dado@W_2.fw.ps\endcsname{32}
\expandafter\def\csname dado@W_2.ap.nsedge\endcsname{0,08}
\expandafter\def\csname dado@W_2.ap.nsedge.cpu\endcsname{0,85}
\expandafter\def\csname dado@W_2.ap.threads\endcsname{11}
\expandafter\def\csname dado@W_2.q4\endcsname{}
\expandafter\def\csname dado@W_2.fwest\endcsname{---}
\expandafter\def\csname dado@W_3.V\endcsname{100.000}
\expandafter\def\csname dado@W_3.E\endcsname{6.099.941}
\expandafter\def\csname dado@W_3.deg\endcsname{122,0}
\expandafter\def\csname dado@W_3.neg\endcsname{não}
\expandafter\def\csname dado@W_3.fwmem\endcsname{120\,GB}
\expandafter\def\csname dado@W_3.fwmem.gib\endcsname{111,8\,GiB}
\expandafter\def\csname dado@W_3.bf\endcsname{102,7\,ms}
\expandafter\def\csname dado@W_3.bf.rounds\endcsname{25}
\expandafter\def\csname dado@W_3.bf.runs\endcsname{10}
\expandafter\def\csname dado@W_3.fw\endcsname{---}
\expandafter\def\csname dado@W_3.ap\endcsname{2,9\,h}
\expandafter\def\csname dado@W_3.ap.sample\endcsname{200}
\expandafter\def\csname dado@W_3.ops.fw\endcsname{$1{,}00\cdot10^{15}$}
\expandafter\def\csname dado@W_3.ops.bf\endcsname{$1{,}52\cdot10^{13}$}
\expandafter\def\csname dado@W_3.ops.ratio\endcsname{65,6$\times$}
\expandafter\def\csname dado@W_3.bf.nsedge\endcsname{0,67}
\expandafter\def\csname dado@W_3.ap.nsedge\endcsname{0,68}
\expandafter\def\csname dado@W_3.ap.nsedge.cpu\endcsname{0,68}
\expandafter\def\csname dado@W_3.ap.threads\endcsname{1}
\expandafter\def\csname dado@W_3.fwest\endcsname{8,9\,h}
\expandafter\def\csname dado@W_4.V\endcsname{1.000.000}
\expandafter\def\csname dado@W_4.E\endcsname{30.999.978}
\expandafter\def\csname dado@W_4.deg\endcsname{62,0}
\expandafter\def\csname dado@W_4.neg\endcsname{não}
\expandafter\def\csname dado@W_4.fwmem\endcsname{12\,TB}
\expandafter\def\csname dado@W_4.fwmem.gib\endcsname{11.175,9\,GiB}
\expandafter\def\csname dado@W_4.bf\endcsname{1,47\,s}
\expandafter\def\csname dado@W_4.bf.rounds\endcsname{36}
\expandafter\def\csname dado@W_4.bf.runs\endcsname{10}
\expandafter\def\csname dado@W_4.fw\endcsname{---}
\expandafter\def\csname dado@W_4.ap\endcsname{14,4\,dias}
\expandafter\def\csname dado@W_4.ap.sample\endcsname{30}
\expandafter\def\csname dado@W_4.ops.fw\endcsname{$1{,}00\cdot10^{18}$}
\expandafter\def\csname dado@W_4.ops.bf\endcsname{$1{,}12\cdot10^{15}$}
\expandafter\def\csname dado@W_4.ops.ratio\endcsname{896$\times$}
\expandafter\def\csname dado@W_4.bf.nsedge\endcsname{1,31}
\expandafter\def\csname dado@W_4.ap.nsedge\endcsname{1,11}
\expandafter\def\csname dado@W_4.ap.nsedge.cpu\endcsname{1,11}
\expandafter\def\csname dado@W_4.ap.threads\endcsname{1}
\expandafter\def\csname dado@W_4.fwest\endcsname{369\,dias}
\expandafter\def\csname dado@W_5.V\endcsname{10.000.000}
\expandafter\def\csname dado@W_5.E\endcsname{39.999.995}
\expandafter\def\csname dado@W_5.deg\endcsname{8,0}
\expandafter\def\csname dado@W_5.neg\endcsname{não}
\expandafter\def\csname dado@W_5.fwmem\endcsname{1,2\,PB}
\expandafter\def\csname dado@W_5.fwmem.gib\endcsname{1.117.587,1\,GiB}
\expandafter\def\csname dado@W_5.bf\endcsname{6,11\,s}
\expandafter\def\csname dado@W_5.bf.rounds\endcsname{42}
\expandafter\def\csname dado@W_5.bf.runs\endcsname{10}
\expandafter\def\csname dado@W_5.fw\endcsname{---}
\expandafter\def\csname dado@W_5.ap\endcsname{2,0\,anos}
\expandafter\def\csname dado@W_5.ap.sample\endcsname{10}
\expandafter\def\csname dado@W_5.ops.fw\endcsname{$1{,}00\cdot10^{21}$}
\expandafter\def\csname dado@W_5.ops.bf\endcsname{$1{,}68\cdot10^{16}$}
\expandafter\def\csname dado@W_5.ops.ratio\endcsname{59.524$\times$}
\expandafter\def\csname dado@W_5.bf.nsedge\endcsname{3,64}
\expandafter\def\csname dado@W_5.ap.nsedge\endcsname{3,84}
\expandafter\def\csname dado@W_5.ap.nsedge.cpu\endcsname{3,84}
\expandafter\def\csname dado@W_5.ap.threads\endcsname{1}
\expandafter\def\csname dado@W_5.fwest\endcsname{1.012\,anos}
\expandafter\def\csname dado@fwref.name\endcsname{\W{2}}
\expandafter\def\csname dado@fwref.fw\endcsname{8,3\,min}
\expandafter\def\csname dado@fwref.fwpm\endcsname{8,3 $\pm$ 0,7\,min}
\expandafter\def\csname dado@neg.list\endcsname{\W{1}}
\expandafter\def\csname dado@W_1.matsrc\endcsname{ciclo negativo}
\expandafter\def\csname dado@W_2.matsrc\endcsname{Floyd-Warshall}
\expandafter\def\csname dado@W_3.matsrc\endcsname{Bellman-Ford, um por destino}
\expandafter\def\csname dado@W_4.matsrc\endcsname{Bellman-Ford, um por destino}
\expandafter\def\csname dado@W_5.matsrc\endcsname{Bellman-Ford, um por destino}
\expandafter\def\csname dado@matrix.bf.list\endcsname{\W{3}, \W{4}, \W{5}}
\expandafter\def\csname dado@W_2.q4.ratio\endcsname{13,1$\times$}
\expandafter\def\csname dado@W_2.q4.winner\endcsname{Bellman-Ford}
\expandafter\def\csname dado@W_2.q4.loser\endcsname{Floyd-Warshall}
\expandafter\def\csname dado@W_3.q4.ratio\endcsname{33,7$\times$}
\expandafter\def\csname dado@W_3.q4.winner\endcsname{Bellman-Ford}
\expandafter\def\csname dado@W_3.q4.loser\endcsname{Floyd-Warshall}
\expandafter\def\csname dado@W_4.q4.ratio\endcsname{282$\times$}
\expandafter\def\csname dado@W_4.q4.winner\endcsname{Bellman-Ford}
\expandafter\def\csname dado@W_4.q4.loser\endcsname{Floyd-Warshall}
\expandafter\def\csname dado@W_5.q4.ratio\endcsname{5.443$\times$}
\expandafter\def\csname dado@W_5.q4.winner\endcsname{Bellman-Ford}
\expandafter\def\csname dado@W_5.q4.loser\endcsname{Floyd-Warshall}
\expandafter\def\csname dado@cmp.name\endcsname{\W{2}}
\expandafter\def\csname dado@cmp.ops\endcsname{31,6$\times$}
\expandafter\def\csname dado@cmp.time\endcsname{13,1$\times$}
\expandafter\def\csname dado@cmp.fw.ps\endcsname{32}
\expandafter\def\csname dado@cmp.bf.ns\endcsname{0,08}
\expandafter\def\csname dado@cmp.cost\endcsname{2,4$\times$ mais}
\expandafter\def\csname dado@w2.check\endcsname{conferem}

\newcommand{\matriz}[1]{\csname matriz@#1\endcsname}

\begin{document}

% ── cabeçalho compacto ──
\noindent\begin{minipage}[c]{2.0cm}
\includegraphics[width=1.7cm]{../comum/minerva.png}
\end{minipage}%
\begin{minipage}[c]{\dimexpr\linewidth-2.0cm\relax}
{\small UFRJ \textperiodcentered{} COS242 Teoria dos Grafos, 2026/2 \textperiodcentered{}
Prof.\ Daniel Ratton Figueiredo}\\[2pt]
{\Large\bfseries Trabalho de Disciplina, Parte 3: caminhos mínimos com pesos negativos}\\[2pt]
{\small Felipe Marroni e Hygor Knust \textperiodcentered{}
\url{https://github.com/knusth/cos-242-grafos}}
\end{minipage}
\par\vspace{2pt}\hrule\vspace{4pt}

\section{Introdução}

A biblioteca agora calcula caminhos mínimos em grafos \emph{direcionados} com pesos que podem ser
negativos. Implementamos o algoritmo de Bellman-Ford, para um destino e para todos os destinos, e o
de Floyd-Warshall, e os aplicamos aos cinco grafos do site da disciplina (Tabela~\ref{tab:grafos}).
O \W{1} tem arestas de peso negativo; os demais não. A última coluna é a memória que o
Floyd-Warshall exige (\S\ref{sec:fw}) contra os \dado{limit} que permitimos a ele.

\begin{table}[!htb]
\centering\small
\caption{Grafos do estudo de caso. Grau médio $=2E/V$ (entrada mais saída); memória do
Floyd-Warshall $=12V^2$ bytes.}
\label{tab:grafos}
\begin{tabular}{@{}lrrrcr@{}}
\toprule
grafo & $V$ & $E$ & grau médio & pesos negativos & memória do Floyd-Warshall \\
\midrule
\tabelaGrafos
\bottomrule
\end{tabular}
\end{table}

\textbf{Validação.} A biblioteca passa em \dado{tests} verificações automáticas
(\cod{tests/graph\_tests.cpp}). Num grafo de 10 mil vértices, a matriz $5\times5$ obtida por
Bellman-Ford, um por destino, é igual célula a célula à do Floyd-Warshall. Conferimos também com
programas independentes em Python: o \W{1} e o \W{2} são fortemente conexos; o \W{1} tem ciclo
negativo (um Bellman-Ford independente, com origem virtual, nunca convergiu); e as distâncias do
\W{2}, calculadas com Dijkstra, \dado{w2.check} com as nossas (Bellman-Ford e matriz $5\times5$).

\section{Decisões de projeto e implementação}

\subsection{Grafos direcionados e representação}

A direção é um parâmetro de execução (\cod{direcionado} ou \cod{nao\_direcionado}), e não de
compilação. A representação é CSR: um grafo direcionado guarda duas listas de adjacência, \cod{forward\_}
(arestas que saem de cada vértice) e \cod{backward\_} (as que chegam), cada uma com $V+1$ deslocamentos e,
por aresta, um vizinho de 4 bytes e um peso de 8. BFS, DFS e Dijkstra percorrem as que saem. Bellman-Ford
calcula a distância de todos os vértices \emph{até} um destino, e por isso percorre as que chegam;
sem direção, as duas listas coincidem e a segunda não é alocada.

\subsection{Bellman-Ford até um destino}

\cod{bellman\_ford(destino)} relaxa as arestas que chegam em cada vértice: se $p\to v$ melhora o caminho
de $p$ ao destino, $\mathit{dist}[p]$ passa a $\mathit{dist}[v]+w(p,v)$ e $\mathit{next}[p]=v$. Os
ponteiros $\mathit{next}$ formam uma árvore com raiz no destino, e o caminho de um vértice sai de seguir
esses ponteiros. Há um só corpo, \cod{relax\_rounds}, com as duas otimizações vistas em aula
(Listagem~\ref{lst:relax}): (1) uma rodada sem mudanças encerra o algoritmo; (2) só os vértices cuja
distância mudou na rodada anterior (\emph{ativos}) têm suas arestas examinadas, sem repetições
graças a um bitmap de marcas. As atualizações valem na própria rodada, então o número de rodadas
depende da ordem das arestas.

\lstinputlisting[float=!htb, caption={Núcleo de uma rodada de \cod{relax\_rounds}
(\cod{include/cos/graph.hpp}, copiado sem alteração).}, label=lst:relax]{relax.txt}

\textbf{Ciclo negativo.} Sem ciclo negativo, um caminho mínimo simples tem até $V-1$ arestas; se a rodada
$V$ ainda muda alguma distância, há um ciclo negativo que alcança o destino. Esse critério, de custo
$O(VE)$, é o último recurso. Três verificações antecipam a resposta: (a) \emph{ciclo entre os
sucessores}: depois de cada rodada, uma busca pelos ponteiros $\mathit{next}$ a partir dos vértices
alterados; como as distâncias só diminuem, todo ciclo entre eles tem peso negativo, e a busca custa $O(V)$
por rodada, com selos que evitam repetir trechos; (b) \emph{melhora do destino}: o destino começa em $0$,
e uma melhora estrita dele (no código, \cod{previous == target}) prova um passeio fechado negativo por ele;
(c) o limite de $V$ rodadas. \cod{has\_negative\_cycle()} usa o mesmo núcleo com um vértice virtual ligado a todos (todos
começam ativos com distância $0$) e por isso acha ciclos em qualquer ponto do grafo,
enquanto \cod{bellman\_ford(destino)} só vê os que alcançam o destino.

\subsection{Floyd-Warshall}
\label{sec:fw}

\textbf{Memória.} Usamos uma única matriz $V\times V$ de distâncias, atualizada no lugar, em vez de uma
por valor de $k$. Ler valores já atualizados é seguro: na iteração $k$, a linha $k$ e a coluna $k$ não
mudam, porque $d[i][k]+d[k][k]=d[i][k]$ quando $d[k][k]=0$. Com a matriz de pais, são $8V^2+4V^2=12V^2$
bytes. A linha $i$ da matriz de pais é a árvore de caminhos mínimos com raiz em $i$: quando $k$ melhora
$d[i][j]$, $\mathit{pai}(i,j)$ passa a ser $\mathit{pai}(k,j)$.

\textbf{Blocos e threads.} As matrizes são divididas em blocos de $64\times64$ (48\,KiB), que ficam na cache
durante 64 valores de $k$. Para cada bloco diagonal, três fases: (1) o bloco diagonal; (2) a linha e a
coluna de blocos dele; (3) os demais blocos, que só leem a linha e a coluna prontas. As threads
reivindicam tarefas por contadores atômicos, com uma barreira entre as fases; as tarefas de uma
fase escrevem em blocos disjuntos, então o resultado não depende do número de threads. O laço
interno (\cod{relax\_segment}) não tem desvios e é vetorizado pelo compilador. Uma diagonal
que cai abaixo de $-$tolerância encerra todas as threads na barreira seguinte: é a saída antecipada
para ciclo negativo.

\textbf{Pais em ciclos de peso zero.} Com a ordem em blocos, empates entre caminhos de mesmo peso
podem ser resolvidos em sentidos opostos, e os pais de uma linha formam um ciclo, em vez de uma
árvore (as distâncias seguem corretas). Uma passagem final detecta esses ciclos e refaz a linha
com uma busca em largura pelas arestas justas.

\textbf{Tolerância numérica.} Os pesos são \cod{double}, e somas como $0{,}1+0{,}2-0{,}3$ não dão $0$
exato: ciclos de soma decimal zero produziam ciclos negativos falsos. Uma relaxação só vale se
melhora a distância em mais que $8\,\varepsilon\,V\max|w|$ (heurística: um caminho tem até $V-1$ somas,
cada uma com erro de $\varepsilon/2$ vezes seu valor). Dijkstra, sem pesos negativos, não usa a margem.

\subsection{Todos os pares por Bellman-Ford}

\cod{all\_pairs\_bellman\_ford(threads)} executa Bellman-Ford uma vez por destino e escreve a linha
de cada destino, contígua, direto na matriz de resultado ($12V^2$ bytes, como no Floyd-Warshall).
Cada thread tem seu espaço de trabalho, reaproveitado entre execuções, e reivindica blocos de
destinos por um contador atômico. Se uma execução acha ciclo negativo, um sinalizador atômico
encerra as demais; a união dos $V$ destinos cobre todos os ciclos negativos. A saída é um tipo à
parte, porque os sucessores (árvore com raiz no \emph{destino}) não são os pais do Floyd-Warshall
(árvore com raiz na \emph{origem}).

\section{Metodologia}

Rodamos tudo num \dado{cpu} com \dado{ram} de memória e \dado{threads} núcleos, compilando em modo
\emph{Release} com \dado{compiler}. Cada algoritmo foi executado 10 vezes por grafo; a tabela traz
a média e o desvio-padrão. Medimos só a execução do algoritmo, sem a leitura do arquivo.
O Bellman-Ford até um destino roda em uma thread; o Floyd-Warshall e o todos-pares completo, em
\dado{threads} threads. O Floyd-Warshall só roda se $12V^2$ bytes couberem em \dado{limit}.
Nos grafos em que o todos-pares por Bellman-Ford levaria horas ou dias, medimos $K$ destinos
igualmente espaçados ($K=\dado{W_3.ap.sample}$ no \W{3}, \dado{W_4.ap.sample} no \W{4},
\dado{W_5.ap.sample} no \W{5}), em uma thread, e multiplicamos a média por $V$; esses valores vêm em
\emph{itálico}. Valores com \textsuperscript{$\dagger$} são do \W{1} (veja a Seção~\ref{sec:q3}).
Medimos o \W{3}, o \W{4} e o \W{5} duas vezes: na primeira passada, as execuções de uma thread
ficaram cerca de $2\times$ mais lentas, com pouca variação, o que é compatível com o macOS tê-las
agendado nos núcleos de eficiência (os arquivos ficam em \cod{bench/results/parte3/primeira\_execucao}); o
relatório usa a segunda passada, com a máquina ociosa.

\section{Estudos de caso}

\subsection{Distâncias por Bellman-Ford até o vértice 100}

\begin{table}[!htb]
\centering\small
\caption{Distância dos vértices 10, 20 e 30 ao vértice 100 (Bellman-Ford). No \W{1}, há um ciclo
negativo e a distância não está definida: vale $-\infty$.}
\label{tab:dist}
\begin{tabular}{@{}r*{5}{r}@{}}
\toprule
origem & \W{1} & \W{2} & \W{3} & \W{4} & \W{5} \\
\midrule
\tabelaDistancias
\bottomrule
\end{tabular}
\end{table}

No \W{1}, o Bellman-Ford detectou o ciclo negativo em \dado{W_1.bf.rounds} rodadas, e a consulta no
grafo inteiro (\cod{has\_}\allowbreak\cod{negative\_cycle}) também o encontra. Como o \W{1} é fortemente conexo,
o ciclo alcança o destino e é alcançado por 10, 20 e 30, então $d(s,100)=-\infty$ para toda
origem $s$: nenhum número seria correto. Nesse grafo, a diferença entre as duas consultas não
importa; ela existiria num grafo que não fosse fortemente conexo. Os caminhos mínimos dos demais
grafos estão no relatório gerado por \cod{cos\_part3} (\url{bench/results/parte3}).

\subsection{Distâncias por Floyd-Warshall entre 10, 20, 30, 40 e 50}

No \W{1}, as 25 distâncias também valem $-\infty$ (o Floyd-Warshall detectou o ciclo pela
diagonal negativa e parou). No \W{2}, a matriz vem do Floyd-Warshall (Figura~\ref{fig:matrizes}). Do \W{3}
ao \W{5}, ele não cabe na memória (Tabela~\ref{tab:grafos}); as matrizes desses grafos vêm do
Bellman-Ford, um por destino, e não do Floyd-Warshall.

\begin{figure}[!htb]
\centering\small\setlength{\tabcolsep}{4.5pt}
\begin{minipage}[t]{0.48\linewidth}\centering
\matriz{W_2}\\[1pt]\W{2}, \dado{W_2.matsrc}
\end{minipage}\hfill
\begin{minipage}[t]{0.48\linewidth}\centering
\matriz{W_3}\\[1pt]\W{3}, \dado{W_3.matsrc}
\end{minipage}\\[5pt]
\begin{minipage}[t]{0.48\linewidth}\centering
\matriz{W_4}\\[1pt]\W{4}, \dado{W_4.matsrc}
\end{minipage}\hfill
\begin{minipage}[t]{0.48\linewidth}\centering
\matriz{W_5}\\[1pt]\W{5}, \dado{W_5.matsrc}
\end{minipage}
\caption{Distância da origem (linha) ao destino (coluna) entre os vértices 10, 20, 30, 40 e 50.
Os grafos são direcionados e a matriz não é simétrica.}
\label{fig:matrizes}
\end{figure}

\subsection{Tempo médio de cada algoritmo}
\label{sec:q3}

\begin{table}[!htb]
\centering\small
\caption{Tempo médio $\pm$ desvio-padrão de 10 execuções, sem leitura do arquivo. Bellman-Ford até o
vértice 100: uma thread; rodadas executadas. Floyd-Warshall e todos-pares por Bellman-Ford (BF):
\dado{threads} threads, exceto os valores em itálico, que são estimados com $K$ destinos em uma
thread. ``memória'': não cabe em \dado{limit}. $^\dagger$: o \W{1} tem ciclo negativo e o
algoritmo para ao detectá-lo; o tempo é de detecção, não de uma execução completa.}
\label{tab:tempos}
\setlength{\tabcolsep}{6pt}
\begin{tabular}{@{}lrrrr@{}}
\toprule
grafo & rodadas & BF até o vértice 100 & Floyd-Warshall & BF todos-pares \\
\midrule
\tabelaTempos
\bottomrule
\end{tabular}
\end{table}

\textbf{\W{1}.} Os três tempos são de detecção do ciclo negativo: o Bellman-Ford para depois de
\dado{W_1.bf.rounds} rodadas, em \dado{W_1.bf}, e o Floyd-Warshall, ao achar uma diagonal
negativa, em \dado{W_1.fw}. O todos-pares para na primeira
busca por destino que acha o ciclo (\dado{W_1.ap}). Não são comparáveis aos tempos dos outros grafos, nem
entre si, pois cada algoritmo interrompe em um ponto diferente.

\textbf{Bellman-Ford até um destino.} Num grafo sem ciclo negativo, o custo é o das arestas que chegam
nos vértices ativos, no máximo $R\cdot E$ para $R$ rodadas. No \W{3}, \dado{W_3.bf.rounds} rodadas sobre
\dado{W_3.E} arestas levam \dado{W_3.bf}, no máximo \dado{W_3.bf.nsedge}\,ns por aresta examinada (limite
superior: nem toda aresta é examinada em toda rodada). O \W{2} leva \dado{W_2.bf} em \dado{W_2.bf.rounds}
rodadas; o \W{4}, \dado{W_4.bf} em \dado{W_4.bf.rounds}; o \W{5}, \dado{W_5.bf} em \dado{W_5.bf.rounds}.

\textbf{Floyd-Warshall.} O \W{2} levou \dado{W_2.fw} por execução em \dado{threads} threads, ou
\dado{W_2.fw.ps}\,ps por relaxação de $V^3$. Do \W{3} em diante, ele não roda; a Seção~\ref{sec:q4} estima o tempo.

\subsection{Todos os pares: Bellman-Ford contra Floyd-Warshall}
\label{sec:q4}

\begin{table}[!htb]
\centering\small
\caption{Floyd-Warshall (FW) contra Bellman-Ford de todos os pares (BF). Operações: FW faz $V^3$
relaxações; BF faz, no máximo, $V\cdot R\cdot E$ exames de aresta, com $R$ as rodadas até o vértice
100 (limite superior). Em itálico, estimativas: o FW do \W{3} em diante é a média do \W{2} vezes
$(V/V_{\W{2}})^3$; o BF vem de $K$ destinos em uma thread. A razão FW/BF divide o tempo do BF pelo
número de threads do FW (escalonamento linear, nas estimativas); acima de 1, o BF é mais rápido.}
\label{tab:q4}
\setlength{\tabcolsep}{5pt}
\begin{tabular}{@{}lrrrrr@{}}
\toprule
grafo & operações do FW & operações do BF (limite) & tempo do FW & tempo do BF & FW/BF \\
\midrule
\tabelaQuatro
\bottomrule
\end{tabular}
\end{table}

\textbf{Operações.} O Floyd-Warshall faz $V^3$ relaxações, qualquer que seja a densidade:
\dado{W_2.ops.fw} no \W{2}. O Bellman-Ford faz $V$ execuções, e cada uma examina só as arestas que
chegam nos vértices ativos, no máximo $R\cdot E$; com as $R=\dado{W_2.bf.rounds}$ rodadas do destino 100, são
no máximo \dado{W_2.ops.bf} exames, \dado{cmp.ops} menos. Esse limite é folgado (um vértice só fica ativo
quando sua distância muda), então a diferença real de operações é maior que essa. O custo por operação
compensa só em parte: a relaxação do Floyd-Warshall é sequencial e vetorizada (\dado{cmp.fw.ps}\,ps de
tempo de parede, com \dado{threads} threads), e o Bellman-Ford gasta \dado{cmp.bf.ns}\,ns por aresta do
limite, \dado{cmp.cost} por operação. A diferença de operações supera a de custo, e o Bellman-Ford foi
\dado{W_2.q4.ratio} mais rápido no \W{2}: \dado{W_2.ap} contra \dado{W_2.fw}. Com a estimativa por $V^3$, a vantagem
cresce com $V$ (\dado{W_3.q4.ratio} no \W{3}, \dado{W_4.q4.ratio} no \W{4}, \dado{W_5.q4.ratio} no \W{5}),
porque $V^3$ cresce mais depressa que $V\cdot R\cdot E$; a estimativa ignora efeitos de cache.

\textbf{Memória.} O Floyd-Warshall precisa da matriz inteira ($12V^2$ bytes): \dado{W_3.fwmem} no \W{3},
\dado{W_4.fwmem} no \W{4} e \dado{W_5.fwmem} no \W{5}. Esse limite pesa mais que o tempo. A
biblioteca também guarda a matriz de resultado no todos-pares por Bellman-Ford, mas o algoritmo
só precisa de $O(V)$ por execução: se cada linha for consumida ao ser calculada, ele roda em
qualquer um dos grafos, e a medição por amostragem de destinos é a de uma implementação assim.
Com a estimativa por $V^3$, o Floyd-Warshall levaria \dado{W_3.fwest} no \W{3},
\dado{W_4.fwest} no \W{4} e \dado{W_5.fwest} no \W{5}; o Bellman-Ford, em uma thread,
\dado{W_3.ap} no \W{3}, \dado{W_4.ap} no \W{4} e \dado{W_5.ap} no \W{5}.
No \W{1} não há comparação: os dois param ao detectar o ciclo negativo.

\section{Conclusão}

No \W{1}, o ciclo negativo torna todas as distâncias indefinidas, e os dois algoritmos o detectam
rapidamente. Nos grafos sem pesos negativos, o Floyd-Warshall só foi possível até o \W{2}; do \W{3}
em diante, sua matriz de $12V^2$ bytes não cabe na memória. No \W{2}, o \dado{W_2.q4.winner} foi
\dado{W_2.q4.ratio} mais rápido que o \dado{W_2.q4.loser} para todos os pares: o Floyd-Warshall faz
$V^3$ relaxações e o Bellman-Ford, no máximo \dado{cmp.ops} menos exames de aresta, e o menor custo
por operação do primeiro não cobre essa diferença. A vantagem estimada cresce com $V$.

\end{document}
